% paper-tle-divergence-atlas -- manuscript
%
% Reproducibility: every figure is regenerated by src/scripts/fig_*.py.
% The underlying sweep is described in sweep/run_sweep.py and archived
% on Zenodo at the DOI listed in the showyourwork.yml `datasets` block.

\documentclass[10pt, letterpaper]{article}
\usepackage[utf8]{inputenc}
\usepackage[T1]{fontenc}
\usepackage[margin=1in]{geometry}
\usepackage{amsmath,amssymb,amsfonts}
\usepackage{graphicx}
\usepackage{booktabs}
\usepackage{natbib}
\usepackage{hyperref}
\usepackage{showyourwork}

\title{How long can you trust a Starlink TLE? \\
       An empirical comparison of SGP4 and high-fidelity propagation
       against operator-updated truth across a megaconstellation}

\author{Dimitrije Jankovic \\
        \small Independent researcher --- astro-tools \\
        \small \href{https://github.com/astro-tools/paper-tle-divergence-atlas}{github.com/astro-tools/paper-tle-divergence-atlas}}

\date{\today}

\begin{document}

\maketitle

\begin{abstract}
\noindent
[Abstract placeholder. Will be written after Results are locked.
Roughly 200 words covering: motivation (TLE ubiquity, megaconstellation scale),
methodology (consecutive-TLE truth construction across $\sim$3000 Starlink
satellite pairs), key findings (TBD --- H1 staleness curves, H2 SGP4 vs.
high-fidelity comparison, H3 solar modulation), and practical implications
for downstream TLE consumers.]
\end{abstract}

\section{Introduction}
\label{sec:introduction}

[TODO: motivation. The growth of public-TLE consumption by smallsat operators,
academic conjunction-analysis pipelines, and amateur trackers; the lack of a
constellation-scale, time-resolved, open-data assessment of how SGP4
propagation accuracy decays with TLE age in modern megaconstellations.]

\section{Background and related work}
\label{sec:background}

[TODO: SGP4 origins \citep{hoots2004}; Vallado's TLE accuracy methodology
\citep{vallado2006, vallado2012}; Crawford on per-satellite divergence
\citep{crawford2007}; megaconstellation-era tracking residuals from LeoLabs
and COMSPOC public statements.]

\section{Data and methods}
\label{sec:methods}

[TODO: CelesTrak archive, sat sampling, pair construction, maneuver filter,
GMAT force-model specification, error metrics.]

\subsection{Spacecraft properties}
\label{sec:methods-spacecraft-props}

Per-satellite dry mass is taken from McDowell's General Catalog of
Artificial Space Objects \citep{mcdowell2020}, which resolves at the
launch-batch level across the Starlink generations in our corpus.
Generation (v1.0 / v1.5 / v2-mini) is classified from the same catalog's
\texttt{PLName} field.

Attitude-averaged effective drag cross-sections are anchored to
\citet{baruah2024}, whose analysis of the February 2022 Starlink loss
event uses ram areas of 1.00\,m$^2$ (collision-avoidance "open book"
attitude) and 4.48\,m$^2$ (operational "shark fin" attitude) at $C_D = 1$
for the v1.5 (Group 4) satellites involved. We adopt the shark-fin value
since station-keeping dominates the duty cycle, rescaled to our
$C_D = 2.2$ convention so $C_D \cdot A$ matches: 2.0\,m$^2$ for v1.x.
For v2-mini, where no equivalent ballistic-coefficient analysis has been
published, we scale by the bus-size ratio reported in
\citet{krebs_starlink_v2mini}, yielding 4.5\,m$^2$ at $C_D = 2.2$; this
is the largest remaining modelling assumption in the spacecraft-properties
layer and is revisited in Section~\ref{sec:discussion}.

\section{Results}
\label{sec:results}

[TODO: prose. H1 addressed in Figures~\ref{fig:sgp4-growth}
and~\ref{fig:hifi-growth}; H2 in Figure~\ref{fig:propagator-scatter}.
Captions below are placeholders pending final Results lock.]

\begin{figure}[ht]
  \centering
  \includegraphics[width=\linewidth]{figures/fig_constellation_map.pdf}
  \caption{Sampled corpus. Pairwise scatter of semi-major axis,
  eccentricity, and inclination for the $\sim$500 Starlink satellites in
  the analysis sample. Altitude shell encoded by color, generation
  encoded by marker shape; diagonal panels show marginal histograms.}
  \label{fig:constellation-map}
  \script{fig_constellation_map.py}
\end{figure}

\begin{figure}[ht]
  \centering
  \includegraphics[width=\linewidth]{figures/fig_sgp4_growth.pdf}
  \caption{SGP4 propagation error vs. $\Delta t$ since epoch, by altitude
  shell within pooled-generation panels. Light scatter shows individual
  $(\Delta t, |\Delta\mathbf{r}|)$ pairs; bold markers connected by lines
  give the per-bucket median and the shaded band gives the 25th--75th
  percentile range. Log-log axes; identical $y$-limits to
  Figure~\ref{fig:hifi-growth}.}
  \label{fig:sgp4-growth}
  \script{fig_sgp4_growth.py}
\end{figure}

\begin{figure}[ht]
  \centering
  \includegraphics[width=\linewidth]{figures/fig_hifi_growth.pdf}
  \caption{High-fidelity propagation error vs. $\Delta t$ since epoch,
  rendered with identical layout and axes to Figure~\ref{fig:sgp4-growth}
  for direct visual comparison.}
  \label{fig:hifi-growth}
  \script{fig_hifi_growth.py}
\end{figure}

\begin{figure}[ht]
  \centering
  \includegraphics[width=\linewidth]{figures/fig_propagator_scatter.pdf}
  \caption{SGP4 vs.\ high-fidelity error per pair, at fixed $\Delta t$
  buckets (rows) and altitude shells (columns). Color encodes pooled
  generation; the dashed $y = x$ line marks parity, so points below are
  pairs where the high-fidelity model beat SGP4.}
  \label{fig:propagator-scatter}
  \script{fig_propagator_scatter.py}
\end{figure}

\begin{figure}[ht]
  \centering
  \includegraphics[width=\linewidth]{figures/fig_powerlaw_fits.pdf}
  \caption{Power-law fits $\|\Delta\mathbf{r}\| \approx A \cdot \Delta t^{k}$
  per (altitude shell $\times$ pooled generation) cell. Bars show the
  point estimate; vertical lines give 95\% bootstrap CIs from 1000
  resamples at the satellite level. Left subpanel per row: coefficient
  $A$ on a log axis; right subpanel: exponent $k$ with $k = 1$ marked.
  Numerical values are listed in Table~\ref{tab:powerlaw}.}
  \label{fig:powerlaw-fits}
  \script{fig_powerlaw_fits.py}
\end{figure}

\begin{table}[ht]
  \centering
  \caption{Power-law fit parameters per (altitude shell $\times$ pooled
  generation), for SGP4 and high-fidelity propagators. Values are point
  estimate followed by 95\% bootstrap CI.}
  \label{tab:powerlaw}
  % Auto-generated by src/scripts/fig_powerlaw_fits.py — do not edit.
\begin{tabular}{llrrrrrrrr}
\toprule
 & & \multicolumn{4}{c}{SGP4} & \multicolumn{4}{c}{high-fid} \\
\cmidrule(lr){3-6} \cmidrule(lr){7-10}
shell & generation & $A$ [95\% CI] & $k$ [95\% CI] & $R^{2}$ & $p_{k=1}/p_{k=2}$ & $A$ [95\% CI] & $k$ [95\% CI] & $R^{2}$ & $p_{k=1}/p_{k=2}$ \\
\midrule
540 km & v1.x & 0.19 [0.17, 0.2] & 0.84 [0.82, 0.87] & 0.40 & $<\!10^{-3}$ / $<\!10^{-3}$ & \textbf{0.13 [0.12, 0.14]} & 1.3 [1.3, 1.3] & 0.68 & $<\!10^{-3}$ / $<\!10^{-3}$ \\
550 km & v1.x & 0.23 [0.19, 0.28] & 0.75 [0.69, 0.83] & 0.33 & $<\!10^{-3}$ / $<\!10^{-3}$ & \textbf{0.23 [0.19, 0.27]} & 0.99 [0.95, 1] & 0.55 & 0.760 / $<\!10^{-3}$ \\
550 km & v2-mini & \textbf{0.049 [0.045, 0.053]} & 1.4 [1.4, 1.5] & 0.77 & $<\!10^{-3}$ / $<\!10^{-3}$ & 0.15 [0.13, 0.18] & 1.1 [1.1, 1.2] & 0.67 & $<\!10^{-3}$ / $<\!10^{-3}$ \\
560 km & v1.x & \textbf{0.15 [0.13, 0.17]} & 0.86 [0.82, 0.89] & 0.41 & $<\!10^{-3}$ / $<\!10^{-3}$ & 0.15 [0.13, 0.17] & 1.1 [1.1, 1.2] & 0.58 & $<\!10^{-3}$ / $<\!10^{-3}$ \\
560 km & v2-mini & \textbf{0.073 [0.066, 0.083]} & 1.3 [1.3, 1.4] & 0.73 & $<\!10^{-3}$ / $<\!10^{-3}$ & 0.18 [0.16, 0.2] & 1.2 [1.1, 1.2] & 0.66 & $<\!10^{-3}$ / $<\!10^{-3}$ \\
\bottomrule
\end{tabular}

\end{table}

\begin{figure}[ht]
  \centering
  \includegraphics[width=\linewidth]{figures/fig_error_decomposition.pdf}
  \caption{Error decomposition into along-track, cross-track, and radial
  components of $\Delta\mathbf{r}$ in the truth state's RSW frame.
  Per-bucket medians with IQR ribbons; solid lines are SGP4, dashed
  lines are high-fidelity; color encodes pooled generation. Confirms
  the along-track-dominant character of long-$\Delta t$ error.}
  \label{fig:error-decomposition}
  \script{fig_error_decomposition.py}
\end{figure}

\begin{figure}[ht]
  \centering
  \includegraphics[width=\linewidth]{figures/fig_solar_modulation.pdf}
  \caption{Per-satellite SGP4 staleness coefficient $A$ vs.\ the mean
  solar F10.7 flux over the satellite's corpus time window, restricted
  to the drag-dominated sub-550\,km shell. Markers colored by pooled
  generation; thick lines are per-generation rolling medians.
  Addresses H3: error growth under drag-driven conditions should track
  solar EUV proxy.}
  \label{fig:solar-modulation}
  \script{fig_solar_modulation.py}
\end{figure}

\section{Discussion}
\label{sec:discussion}

[TODO: practical implications for downstream TLE consumers; whether
investing in high-fidelity OD is justified for public TLE workflows;
limitations.]

\section{Conclusion}
\label{sec:conclusion}

[TODO.]

\appendix

\section{Maneuver filter calibration}
\label{app:maneuver-filter}

The pair-construction pipeline drops consecutive TLE pairs whose
mean-motion derived semi-major-axis difference $|\Delta a|$ exceeds a
fixed threshold, on the grounds that such jumps mark station-keeping
burns and so violate the no-maneuver assumption behind the truth
construction. Figure~\ref{fig:maneuver-filter} shows the
$|\Delta a|$ distribution across all consecutive pairs in the raw
cache, with the 100\,m threshold marked.

\begin{figure}[ht]
  \centering
  \includegraphics[width=\linewidth]{figures/fig_maneuver_filter.pdf}
  \caption{Distribution of $|\Delta a|$ between consecutive Starlink
  TLEs across the raw cache. The dashed line marks the 100\,m maneuver
  filter threshold; the two modes (quiet vs.\ maneuvering pairs) sit
  cleanly on either side.}
  \label{fig:maneuver-filter}
  \script{fig_maneuver_filter.py}
\end{figure}

\section*{Data and code availability}
\label{sec:availability}

All code is available at
\url{https://github.com/astro-tools/paper-tle-divergence-atlas}
under the MIT license. The full sweep output bundle is archived on Zenodo
[DOI TBD at v0.1.0 release]. Input TLE data is publicly available from
\href{https://celestrak.org}{CelesTrak}.

\bibliographystyle{plainnat}
\bibliography{bib}

\end{document}
